\subsection*{Chapter 9: Optimizers and the Training Loop from First Principles}

\paragraph{Exercise 9.1.} Without correction, $m_1=(1-\beta_1)g$ and $v_1=(1-\beta_2)g^2$, so
the step is $\eta\,(1-\beta_1)|g|/\sqrt{1-\beta_2}\,|g| = \eta\cdot0.1/0.0316\approx3.16\,\eta$
per coordinate. With $\beta_2=0.95$ it is $0.1/\sqrt{0.05}\approx0.45\,\eta$. The second-moment
correction matters more because $\sqrt{1-\beta_2^t}$ stays far below $1-\beta_1^t$ for many
steps when $\beta_2$ is close to 1. The uncorrected ratio $(1-\beta_1^t)/\sqrt{1-\beta_2^t}$
peaks at about $6.6$ around $t=12$ for $(0.9, 0.999)$, so an uncorrected Adam takes steps
several times larger than $\eta$ throughout the first hundreds of steps.

\paragraph{Exercise 9.2.} Adam + L2: the gradient is $\lambda\theta$, so $\hat m\approx\lambda\theta$
and $\sqrt{\hat v}\approx\lambda|\theta|$ once the moments have adapted, giving a step of
$\eta\,\mathrm{sign}(\theta)$: a \emph{linear} shrink of $\eta$ per step regardless of $\lambda$,
until $\theta$ reaches $\approx0$ and oscillates at amplitude $\sim\eta$. AdamW: the adaptive
term is zero ($g=0$, so $m=0$) and $\theta_t=(1-\eta\lambda)^t\theta_0$ decays geometrically with
rate set by $\lambda$. The coupled version's decay strength therefore depends on the gradient
statistics, which is the problem decoupling fixes.

\paragraph{Exercise 9.3.} (a) SGD with momentum in fp32: read $p,g,b$ (12\,B), write $p,b$
(8\,B), 20\,B/param, 5.97\,ms per billion parameters. (b) Lion with fp32 master weights: read bf16
$g$ (2), master (4), $m$ (4); write master (4), $m$ (4), bf16 copy (2): 20\,B/param, 5.97\,ms/B.
(c) AdamW with bf16 $m,v$ and fp32 master: read $g$ (2), master (4), $m$ (2), $v$ (2); write
master (4), $m$ (2), $v$ (2), bf16 copy (2): 20\,B/param, 5.97\,ms/B. For comparison, fp32 or
standard mixed-precision AdamW is 28\,B/param, 8.36\,ms/B.

\paragraph{Exercise 9.4.} $310\times10\times5\,\mu$s $=15.5$\,ms of launch overhead per step,
against a memory-bound time of $28\times0.596\times10^9/3.35\times10^{12}\approx5.0$\,ms. The loop
implementation is launch-bound, three times slower than the bandwidth limit, before counting the
kernels' own runtime. The \code{foreach} path issues about ten launches per parameter group
($\sim$50\,$\mu$s); a fused kernel approaches the 5\,ms bound.

\paragraph{Exercise 9.5.} Qwen3-0.6B: $N\approx0.596\times10^9$, so model state is $16N\approx9.5$\,GB.
At $T=4096$ in bf16 the saved tensors of \cref{sec:optim:account} give (per sequence):
(i) naive attention, $\approx1.26$\,GB per layer, dominated by the $2HT^2$ score and probability
tensors (1.07\,GB), so $28\times1.26\approx35.2$\,GB, plus about 5\,GB for fp32 logits and softmax
($T\cdot V\cdot4\cdot2$ bytes with $V=151{,}936$). That is $\approx40$\,GB per sequence:
$\text{mem}(B)\approx40.2B+9.5$\,GB, which does not fit even $B=1$ on a 24\,GB GPU. (ii) With
FlashAttention, $\approx0.185$\,GB per layer, so $\approx5.2+5.0=10.2$\,GB per sequence:
$\text{mem}(B)\approx10.2B+9.5$\,GB, and $B=1$ fits. The logits are now half the activation memory:
computing the loss in chunks (\cref{ch:distill}) would allow $B=2$, and activation
checkpointing (\cref{sec:pretrain:efficiency}) more.

\paragraph{Exercise 9.6.} Their schedule restarts at step 0: if it has warmup, the learning rate
ramps up again from near zero to the peak, then follows a complete cosine compressed into the
remaining steps. The original run was somewhere down its cosine (say at $0.4\,\eta_{\max}$),
so the resumed run jumps back to $\eta_{\max}$. The loss curve shows a bump or spike shortly
after resumption (a ``warm restart''), then decays again. The final loss is usually worse than the
uninterrupted run's, and the run is not reproducible. The fix is to construct the schedule with the
original \code{total} and call it with the restored global step.

\paragraph{Exercise 9.7.} The global norm is $\sqrt{3^2+4^2}=5$, so both ranks must scale by
$1/5$; the result has norm 1 and the original direction. Independent clipping scales rank 0 by
$1/3$ and rank 1 by $1/4$: each shard has norm 1, the combined norm is $\sqrt2$, and the
direction changes (the components' ratio goes from $3:4$ to $1:1$). The fix is one scalar
collective: each rank computes $s_r=\|g_r\|^2$, then \code{all\_reduce(SUM)} over the
tensor-parallel group, and every rank uses $\sqrt{\sum_r s_r}$. Replicated parameters, such as
norm gains that every rank holds, must be counted once, not once per rank.

\paragraph{Exercise 9.8.} With $X=U\Sigma V^\top$, $XX^\top=U\Sigma^2U^\top$, hence
$(XX^\top)X=U\Sigma^3V^\top$ and $(XX^\top)^2X=U\Sigma^5V^\top$. So
$p(X)=U(a\Sigma+b\Sigma^3+c\Sigma^5)V^\top=U\,p(\Sigma)\,V^\top$, and singular vectors are
preserved. The scalar map only pushes values in $(0,1]$ into the band near 1. For $\sigma$ much
above 1 the $c\sigma^5$ term dominates and iterating diverges, so the input is divided by its
Frobenius norm, which upper-bounds $\sigma_{\max}$. Muon is restricted to 2-D hidden matrices
because 1-D gains and biases have no singular structure, and because embedding and unembedding
matrices are lookups or classifiers over tokens: each row belongs to one token and receives
gradient only when that token appears, and orthogonalising would redistribute update magnitude
across rare and frequent tokens. The original Muon recipe keeps them on AdamW.

\paragraph{Exercise 9.9.} Number the $n$ valid windows (e.g.\ disjoint windows $i\cdot T$).
Each epoch draws $a$ with $\gcd(a,n)=1$ and an offset $c$; the $k$-th sample is window
$(a k + c)\bmod n$, a bijection on $\{0,\dots,n-1\}$ computed in $O(1)$ without storing a
permutation. The state is (epoch, $k$, $a$, $c$), or the seed from which $a$ and $c$ are derived.
Affine permutations mix poorly (consecutive samples are an arithmetic progression of windows), so
production samplers use a small keyed Feistel network or a block-shuffled order instead; the
state has the same shape.

\paragraph{Exercise 9.10.} Expect the following pattern; exact values depend on your hardware
and seeds. With clipping, the largest learning rates spike but recover; our run at $\eta=1$
spiked from 5.8 to 7.3 and then trained. With clipping disabled, the loss at the highest
learning rates becomes non-finite within tens of steps, which defines the divergence point
$\eta_{\text{div}}$. The best learning rate typically sits 3--30$\times$ below $\eta_{\text{div}}$
on a broad plateau. Removing QK-norm lets attention logits grow with the weights, which lowers
$\eta_{\text{div}}$, often by several times, and usually moves the optimum down with it: QK-norm
buys stability headroom (\cref{sec:archdesign:stability}).
